\chapter{Comments on Existence Theory for Parabolic PDE}

\section*{4.1 \quad Linear scalar PDE}

Let $\Omega\subset\mathbb R^n$ and let $u:\Omega\to\mathbb R$. A second-order
equation
\[
\frac{\partial u}{\partial t}
 =a_{ij}\partial_i\partial_j u+b_i\partial_i u+cu,
\tag{4.1.1}
\]
with $\partial_i=\frac{\partial}{\partial x^i}$ and smooth coefficients
$a_{ij},b_i,c:\Omega\to\mathbb R$, is parabolic when there is a constant
$\lambda>0$ such that
\[
a_{ij}\xi_i\xi_j\geq\lambda\lvert\xi\rvert^2
\tag{4.1.2}
\]
for every $\xi\in\mathbb R^n$. In particular, the heat equation
$\frac{\partial u}{\partial t}=\Delta u$ is parabolic.

Let $\M$ be closed and let $L:C^\infty(\M)\to C^\infty(\M)$ be a second-order
operator which, in local coordinates and with smooth local coefficients, has the form
\[
L(u)=a_{ij}\partial_i\partial_j u+b_i\partial_i u+cu.
\]
The equation
\[
\frac{\partial u}{\partial t}=L(u)
\tag{4.1.3}
\]
is parabolic when $a_{ij}(x)$ is positive definite for every $x\in\M$; this
condition is coordinate independent. For every smooth $u_0:\M\to\mathbb R$,
there is a smooth solution on $\M\times[0,\infty)$ satisfying
\[
\begin{cases}
\dfrac{\partial u}{\partial t}=L(u),\\
u(0)=u_0.
\end{cases}
\]
If $u$ and $v$ solve the corresponding equation on $\M\times[0,T]$ and either
$u(0)=v(0)$ or $u(T)=v(T)$, then $u(t)=v(t)$ for every $t\in[0,T]$.

\section*{4.2 \quad The principal symbol}

\begin{definition}
\label{topping-pde-principal-symbol-scalar}
\notready
\source{topping:page-0055}
\dcref{ch4:2.1}
\group{topping-pde-existence}
For the scalar operator $L$ above, its principal symbol is the function
$\sigma(L):T^*\M\to\mathbb R$ given in local coordinates by
\[
\sigma(L)(x,\xi)=a_{ij}(x)\xi_i\xi_j.
\]
Equivalently, if $d\phi(x)=\xi$ for smooth $\phi,f:\M\to\mathbb R$, then
\[
\sigma(L)(x,\xi)f(x)
 =\lim_{s\to\infty}s^{-2}e^{-s\phi(x)}L(e^{s\phi}f)(x).
\tag{4.2.2}
\]
\end{definition}

\begin{remark}
\label{topping-pde-principal-symbol-limit-justification}
\notready
\source{topping:page-0056}
\dcref{ch4:2.2}
\group{topping-pde-existence}
The coordinate formula and the limit formula agree.
\end{remark}
\begin{proof}
\sketch Differentiating $e^{s\phi}f$ shows that the order-$s^2$ term in
$e^{-s\phi(x)}L(e^{s\phi}f)(x)$ is
$s^2a_{ij}\partial_i\phi(x)\partial_j\phi(x)f(x)$. Since
$d\phi(x)=\xi$, division by $s^2$ and passage to the limit gives
$a_{ij}\xi_i\xi_jf(x)$.
\end{proof}

The equation (4.1.3) is parabolic precisely when
$\sigma(L)(x,\xi)>0$ for every $(x,\xi)\in T^*\M$ with $\xi\neq0$.

\begin{example}
\label{topping-pde-laplace-beltrami-example}
\notready
\source{topping:page-0056}
\dcref{ch4:2.2}
\group{topping-pde-existence}
For the Laplace--Beltrami operator,
\[
\Delta=\frac{1}{\sqrt g}\partial_i\bigl(\sqrt g\,g^{ij}\partial_j\bigr)
 =g^{ij}\partial_i\partial_j+\text{lower order terms},
\qquad g=\det(g_{ij}),
\]
one has
\[
\sigma(\Delta)(x,\xi)=g^{ij}\xi_i\xi_j=\lvert\xi\rvert^2>0.
\]
Thus the heat equation defined by this operator is parabolic.
\end{example}

\section*{4.3 \quad Generalisation to vector bundles}

Let $E$ be a smooth vector bundle over a closed manifold $\M$, and let
$v\in\Gamma(E)$. In a local frame $\{e_\alpha\}$, a linear second-order operator
$L:\Gamma(E)\to\Gamma(E)$ has the form
\[
L(v)=\bigl[a^{ij}_{\alpha\beta}\partial_i\partial_jv^\beta
 +b^i_{\alpha\beta}\partial_iv^\beta+c_{\alpha\beta}v^\beta\bigr]e_\alpha.
\]
For
\[
\frac{\partial v}{\partial t}=L(v),
\tag{4.3.1}
\]
the principal symbol is the homomorphism
$\sigma(L):\Pi^*(E)\to\Pi^*(E)$, where $\Pi:T^*\M\to\M$ is the bundle
projection and the fibre over $(x,\xi)$ is $E_x$, defined by
\[
\sigma(L)(x,\xi)v
 =\bigl(a^{ij}_{\alpha\beta}\xi_i\xi_jv^\beta\bigr)e_\alpha.
\tag{4.3.2}
\]
Equivalently, for $v\in\Gamma(E)$ and $\phi:\M\to\mathbb R$ with
$d\phi(x)=\xi$,
\[
\sigma(L)(x,\xi)v
 =\lim_{s\to\infty}s^{-2}e^{-s\phi(x)}L(e^{s\phi}v)(x).
\tag{4.3.3}
\]

If $\Delta$ is the connection Laplacian of a vector bundle with connection,
then
\[
\Delta(e^{s\phi}v)
 =e^{s\phi}\bigl(s^2\lvert d\phi\rvert^2v+\text{terms of lower order in }s\bigr),
\]
so $\sigma(\Delta)(x,\xi)=\lvert\xi\rvert^2\id$.

The vector-bundle equation (4.3.1) is strictly parabolic (also called strongly
parabolic) if there is $\lambda>0$ such that
\[
\langle\sigma(L)(x,\xi)v,v\rangle
 \geq\lambda\lvert\xi\rvert^2\lvert v\rvert^2
\tag{4.3.4}
\]
for every $(x,\xi)\in T^*\M$ and $v\in\Gamma(E)$.

\begin{remark}
\label{topping-pde-vector-parabolic-fibre-metric}
\notready
\source{topping:page-0058}
\dcref{ch4:3.1}
\group{topping-pde-existence}
Equation (4.3.1) is parabolic if (4.3.4) holds for any choice of fibre metric
$\langle\cdot,\cdot\rangle$ on $E$.
\end{remark}

\begin{remark}
\label{topping-pde-vector-parabolic-positive-multiple}
\notready
\source{topping:page-0058}
\dcref{ch4:3.2}
\group{topping-pde-existence}
On a closed manifold, (4.3.1) is parabolic if, for every
$(x,\xi)\in T^*\M$ with $\xi\neq0$, $\sigma(L)(x,\xi)$ is a positive multiple
of the identity on $E_x$.
\end{remark}

Consider the quasilinear equation
\[
\frac{\partial v}{\partial t}=P(v),
\tag{4.3.5}
\]
where $P:\Gamma(E)\to\Gamma(E)$ is given, in local coordinates and a local frame,
\[
P(v)=\bigl[a^{ij}{}_{\alpha\beta}(x,v,\nabla v)\partial_i\partial_jv^\beta
 +b^\alpha(x,v,\nabla v)\bigr]e_\alpha.
\]
The Ricci flow equation $\frac{\partial g}{\partial t}=-2\operatorname{Ric}(g)$
is of this form. For a fixed time-independent $w\in\Gamma(E)$, equation
(4.3.5) is parabolic at $w$ when its linearisation
\[
\frac{\partial v}{\partial t}=[DP(w)]v
\]
is parabolic in the vector-bundle sense.

\section*{4.4 \quad Properties of parabolic equations}

If (4.3.5) is parabolic at $w$ and its coefficients and lower-order terms have
the required regularity, then there are $\varepsilon>0$ and a smooth family
$v(t)\in\Gamma(E)$ for $t\in[0,\varepsilon]$ satisfying
\[
\begin{cases}
\dfrac{\partial v}{\partial t}=P(v),\\
v(0)=w.
\end{cases}
\]
If $v$ and $w$ solve this equation on $[0,\varepsilon]$ and either
$v(0)=w(0)$ or $v(\varepsilon)=w(\varepsilon)$, then
$v(t)=w(t)$ for every $t\in[0,\varepsilon]$.
